\oppgave{7}
Den opprinnelige bitraten og de to roll-off-faktorene er:
\[
    R_{b,1}=6400\enhet{bps},\qquad
    \alpha_1=0.55,\qquad\alpha_2=0.15
\]
Vi antar at kanalbåndbredden beholdes. Båndbredden for den opprinnelige binære overføringen er:
\[
    B=\frac{1+\alpha_1}{2}R_{b,1}
\]
Vi setter inn verdiene:
\[
    B=\frac{1+0.55}{2}\cdot6400=4960\enhet{Hz}
\]
Når roll-off-faktoren reduseres, kan bitraten økes innenfor samme båndbredde:
\[
    R_{b,2}=\frac{2B}{1+\alpha_2}
\]
Den nye bitraten blir:
\[
    R_{b,2}=\frac{2\cdot4960}{1+0.15}
           \approx8626.087\enhet{bps}
\]
Vi kunne også funnet dette direkte fra forholdet mellom roll-off-faktorene:
\[
    R_{b,2}=R_{b,1}\frac{1+\alpha_1}{1+\alpha_2}
           =6400\frac{1.55}{1.15}
           \approx8626.087\enhet{bps}
\]
Økningen i bitrate er:
\[
    \Delta R_b=8626.087-6400\approx2226.087\enhet{bps}
\]
Den prosentvise økningen blir:
\[
    \frac{R_{b,2}-R_{b,1}}{R_{b,1}}\cdot100\%
       =\frac{2226.087}{6400}\cdot100\%
       \approx34.783\%
\]
Informasjonen kan dermed sendes med:
\[
    \boxed{R_{b,2}\approx8.626\enhet{kbps},
    \quad\text{en økning på }34.783\%}
\]
En lavere roll-off-faktor gir mindre overskuddsbåndbredde, men stiller samtidig strengere krav til pulsforming og tidsstyring i mottakeren (mindre robust mot jitter).
